% MT040 / explanation — canonical editable content.
\(x\)是一条待分类记录的特征，\(c\)是候选类别。先验\(P(c)\)表示还没看这条记录时，类别\(c\)有多常见；类条件分布\(p(x\mid c)\)描述已知类别是\(c\)时，什么输入常见；后验\(P(c\mid x)\)表示看过输入\(x\)后，它属于\(c\)类的概率。竖线“\(\mid\)”读作“在已知……的条件下”，条件写在竖线右边。二分类的两个后验概率加起来为1。


Bayes决策规则在错分代价相同时选择后验最大的类。NB通过一个具体假设计算这些后验：给定类别后，各特征独立，因而可以把各维的概率或密度相乘，再结合类别先验。


二分类时，两类的未归一化分数为
\[
s_0=p(x\mid0)P(0),\qquad s_1=p(x\mid1)P(1).
\]


真正的后验还要分别除以共同的\(p(x)=s_0+s_1\)。这个分母是同一个正数，所以只想比较大小时可以省略；先验\(P(c)\)却随类别变化，不能随手删掉。



\begin{itemize}

\item \textbf{A选。} NB先作条件独立的建模假设，再执行上述规则，是一种具体实现。

\item \textbf{B不选。} 补充数值：类条件密度为0.4、0.2，先验为0.1、0.9。两类分数是0.04、0.18，后验约为0.182、0.818，应选类1。单独比较密度会选错。连续量的密度值不是“恰好取这个值的概率”，这里用它比较证据。

\item \textbf{C选，但补上条件。} 独立指\textbf{给定类别之后}独立，即\(p(x\mid c)=\prod_j p(x_j\mid c)\)，其中\(\prod_j\)表示把每个特征\(j\)对应的项相乘。这个假设不要求混合所有类别后的特征也独立。

\item \textbf{D不选。} 决策规则没有规定采用什么概率模型，无法脱离模型谈它必然比NB快。

\item \textbf{E不选。} 三类、十类都可以逐一算后验再取最大值，规则不局限于二分类。

\end{itemize}


\textbf{出处：} 2023A期中 Q1，原卷第2页（5分）；2025A期中 Q1，原卷第2页（5分）


\textbf{答案依据：} 2023A期中Q1随卷答案AC；2025A Q1同题复用。
